\begin{claim}{D.56}{Proved}{The prime sections of a cyclic-six torus action}
Let $a$ have order six in the full once-punctured-torus
loop-outer group of $F=\widehat F(x,y)$, and let
\[
 E=\{(\phi,j)\in\operatorname{Aut}(F)\times C_6:
                                      [\phi]=a^j\}
                                                               \tag{D.56a}
\]
be its centre-free extension, with kernel identified
with $F$ by the inner-automorphism map.
There are exactly three $F$-conjugacy classes of
order-two sections over $\langle a^3\rangle$ and
exactly two $F$-conjugacy classes of order-three
sections over $\langle a^2\rangle$.
Conjugation by an element with quotient label $a^2$
cycles the three former classes; conjugation by an
element with quotient label $a^3$ exchanges the two
latter classes.

Consequently $E$ has exactly one conjugacy class
of subgroups of order two and one of order three,
and these are its only nontrivial finite subgroup
types.  In particular the full extension (D.56a)
does not split, and $a$ has no finite-order
representative in $\operatorname{Aut}(F)$.
For each prime section $S$ of order $p=2$ or $3$,
\[
                 N_E(S)=C_F(S)\times S,
                 \qquad H^1(C_F(S),\F_p)=0.             \tag{D.56b}
\]
The assertions concern these specified extension
and section maps; they do not identify $E$ with
a free profinite product.
\end{claim}
\begin{proof}
Put
\[
 E_2=\pi^{-1}(\langle a^3\rangle),\qquad
 E_3=\pi^{-1}(\langle a^2\rangle),
 \quad \pi:E\longrightarrow\langle a\rangle .          \tag{D.56c}
\]
The involution $a^3$ is positive.
Indeed D.55 embeds $\langle a\rangle$ faithfully
in $\operatorname{GL}_2(\Q_2)$.
An order-six matrix there has its order-three
part in the irreducible quadratic field
$\Q_2(\zeta_3)$; its commuting involution is
the nonidentity scalar $-I$.
For completeness, the centralizer is that field
because the cubic part has irreducible polynomial
$X^2+X+1$, and its only elements squaring to one
are $1,-1$. Faithfulness excludes the former
for $a^3$.
Thus $a^3$ has determinant one at two, and D.43
gives its global positive sign.
D.27 supplies the actual prime-order sections
of both groups in (D.56c).

\emph{The three involution classes and their affine coordinates.}
Choose an order-two section generator $t$ in
$E_2$ and identify $E_2=F\rtimes\langle t\rangle$.
By D.46 and D.55, its action on two-adic
homology is $-I$, and its maximal elementary
abelian two-quotient has coordinates
\[
 (u,\chi_1,\chi_2):E_2\longrightarrow\F_2^3,
 \quad u(t)=1,\quad\chi_1(t)=\chi_2(t)=0,                \tag{D.56d}
\]
where $\chi_1|_F,\chi_2|_F$ form a basis of
$H^1(F,\F_2)$.
Every section generator has $u=1$, so its
coordinate is a point of the affine plane
$\F_2^2$.
The three $F$-conjugacy classes of D.46 give
three distinct such points.
One can read the distinctness directly from
that proof: Brown restriction identifies the
localized degree-zero ring with $\F_2^3$,
and this ring is generated by $\chi_1/u$ and
$\chi_2/u$. Equal values on two section
classes would prevent these two elements
from generating the corresponding two
coordinate factors.

Conjugation by an element $e\in E$ with
$\pi(e)=a^2$ preserves $u$ and acts on the
plane (D.56d) by
\[
                           r\longmapsto Rr+d,           \tag{D.56e}
\]
where the linear part is the reduction of
$A_2(a^2)$ in the chosen homology coordinates.
D.54, at $q=2$, says that $R-I$ is invertible.
Also $R^3=I$, so $R$ is a nonidentity cubic
matrix over $\F_2$ and $I+R+R^2=0$.
The affine map (D.56e) has one fixed point,
and its other three points form a single
orbit of length three: translate its unique
fixed point to zero and use the corresponding
linear statement.
Its invariant three-point set of section
classes must be that orbit.
This proves the asserted cyclic permutation.
The action on the class set factors through
$E/E_2=C_3$, because conjugation by $E_2$
is already $F$-conjugation on its sections,
as proved in D.46.

\emph{The two cubic classes and their affine coordinates.}
Choose an order-three section generator $s$
in $E_3$ with label $a^2$.
Use the actual abelian quotient of D.48 and
the separation proved in D.51:
\[
 (u_3,\chi):E_3\longrightarrow\F_3^2,
 \qquad u_3(s)=1,\quad\chi(s)=0.                        \tag{D.56f}
\]
The two $F$-conjugacy classes of section
generators with prescribed label $u_3=1$
have $\chi$-values $0$ and $b$, with
$b\ne0$.
For an element $e\in E$ with label $a^3$,
conjugation preserves $u_3$ because $E/F$
is abelian.  On the restriction of $\chi$
to $F$ it acts by $-1$, since D.55 gives
$A_3(a^3)=-I$.
The difference between the conjugated
character and $-\chi$ vanishes on $F$,
so is a scalar multiple of $u_3$.
On the line $u_3=1$ the action consequently
has the exact form
\[
                              r\longmapsto-r+d'.        \tag{D.56g}
\]
Preserving the two-point set $\{0,b\}$
forces $d'=b$: the alternative of fixing
both would give $2b=0$ in $\F_3$.
Thus this action exchanges the two classes.
Conjugation by $E_3$ is $F$-conjugation on
its sections, by D.48, so the action factors
through $E/E_3=C_2$.

\emph{The full finite subgroups and the failure of lifting.}
Every finite subgroup of $E$ injects into
$E/F=C_6$, because $F$ is torsion-free.
It is therefore cyclic of order one, two,
three or six.  A subgroup of order six
would normalize its order-two subgroup.
Its element with quotient label $a^2$
would then fix that subgroup's
$F$-conjugacy class, contradicting
(D.56e).  Thus no subgroup of order six
exists.
The two affine actions show respectively
that all order-two subgroups are conjugate
in $E$ and all order-three subgroups are
conjugate in $E$.

For a chosen section $S$ of order two,
the stabilizer of its $F$-conjugacy class
in $E/F$ is exactly $\langle a^3\rangle$;
for order three it is exactly
$\langle a^2\rangle$.
The image of $N_E(S)$ equals this
stabilizer: a conjugator whose class is
fixed can be multiplied by an element
of $F$ to normalize $S$ itself.
Thus the full normalizer is already
its normalizer in $E_p$.
Equations (D.46h) and (D.48g) identify
it as $C_F(S)\times S$ and give the
cohomology assertion in (D.56b).

The first projection of (D.56a) is
injective and identifies $E$ with the
inverse image of $\langle a\rangle$
in $\operatorname{Aut}(F)$.
Indeed the outer class of the first
coordinate determines its label $j$.
A homomorphic lift of $\langle a\rangle$
would give a subgroup of order six in
this exact inverse image, which has
been excluded.
More generally, a finite-order
representative $\phi$ of $a$ would
satisfy $\phi^6\in\operatorname{Inn}(F)$.
The latter group is isomorphic to the
torsion-free group $F$, so
$\phi^6=1$, making $\phi$ a forbidden
element of order six.
This proves the labelled nonlifting
assertion while preserving the actual
prime-order lifts in (D.56c).

The program
\path{code/part_a/torus_order6_sections_d56.py}
checks every affine cubic action on
$\F_2^2$ with nonidentity cubic linear
part, both nonzero cubic section
separations in (D.56f), and the
corresponding quotient stabilizers.
The extension and cohomological
arguments above establish that these
are the actions on the full section
class sets.
\end{proof}
